(a) If $I(Y)$ is generated by linear forms, $Y$ is the intersection of their hyperplanes. Conversely, an ideal generated by linear forms becomes $(x_0,\ldots,x_{q-1})$ after a linear change of coordinates, so it is prime and is the ideal of its nonempty projective zero set.

(b) Choose a basis of $I(Y)_1$ of size $q$. The same coordinate change gives $S(Y)\cong k[x_q,\ldots,x_n]$, so $r=n-q$. Every linear generating set must span $I(Y)_1$, hence its minimal size is $n-r$.

(c) The affine cones are vector subspaces of dimensions $r+1$ and $s+1$, so their intersection has dimension at least $r+s-n+1$. If $r+s\ge n$, it contains a nonzero vector. Whenever the projective intersection is nonempty, it is the projectivization of this vector subspace and is linear of dimension at least $r+s-n$.
